\section{回归实验：效率表现预测}

\subsection{任务设定}

回归目标包括 \texttt{productivity\_score} 与 \texttt{digital\_dependence\_score}。实验使用人口背景、数字行为、生活习惯与工程特征进行预测，但排除 \texttt{high\_risk\_flag}、心理状态结果变量和另一个 outcome 目标变量。

\subsection{模型与指标}

模型包括 Linear Regression、Ridge Regression、Random Forest Regressor 和 Gradient Boosting Regressor。第二阶段补充 5 折交叉验证调参，评价指标包括 MAE、MSE、RMSE 与 $R^2$。

\subsection{\texttt{productivity\_score} 预测}

\placeholdertext{Table placeholder: report/tables/regression\_productivity\_metrics.tex}

\maybefigure{regression_productivity_observed_vs_predicted.png}{效率分数真实值与预测值对比}{fig:regression-prediction}

\maybefigure{regression_productivity_residuals.png}{效率预测残差图}{fig:regression-residuals}

\maybefigure{regression_productivity_permutation_importance.png}{回归任务置换重要性}{fig:regression-importance}

\subsection{\texttt{digital\_dependence\_score} 预测}

\placeholdertext{Table placeholder: report/tables/regression\_digital\_dependence\_metrics.tex}

\maybefigure{regression_digital_dependence_observed_vs_predicted.png}{数字依赖分数真实值与预测值对比}{fig:regression-dependence-prediction}

\maybefigure{regression_digital_dependence_residuals.png}{数字依赖预测残差图}{fig:regression-dependence-residuals}

\maybefigure{regression_digital_dependence_permutation_importance.png}{数字依赖回归任务置换重要性}{fig:regression-dependence-importance}

\subsection{回归目标对比}

\placeholdertext{Table placeholder: report/tables/regression\_target\_comparison.tex}

\maybefigure{regression_target_comparison.png}{两个回归目标的最佳模型表现对比}{fig:regression-target-comparison}
